\documentclass[11pt]{article}
\usepackage[T1]{fontenc}
\usepackage{lmodern}
\usepackage{amsmath,amssymb,amsthm}
\usepackage[margin=1in]{geometry}
\usepackage{microtype}
\usepackage{xurl}
\usepackage[hidelinks,breaklinks]{hyperref}
\hypersetup{pdftitle={Finite-sample density reconstruction from a single causal order},pdfauthor={Daniel Eric Fredriksen},pdfkeywords={partial order, density estimation, copula, finite sample, Lean}}
\newtheorem{theorem}{Theorem}[section]
\newtheorem{lemma}[theorem]{Lemma}
\newtheorem{proposition}[theorem]{Proposition}
\newtheorem{corollary}[theorem]{Corollary}
\theoremstyle{remark}
\newtheorem{remark}[theorem]{Remark}
\newcommand{\R}{\mathbb R}
\newcommand{\K}{\mathcal K}
\newcommand{\TV}{\operatorname{TV}}
\newcommand{\dc}{d_{\mathrm{conf}}}
\newcommand{\rank}{\operatorname{rank}}
\newcommand{\dd}{\,\mathrm d}
\newcommand{\norm}[1]{\lVert #1\rVert_{\infty,D}}
\title{Finite-sample density reconstruction\\from a single causal order}
\author{Daniel Eric Fredriksen\\Quantyra Inc.\\
\url{https://github.com/Quantyra/quantyra-null-cone-lean}}
\date{9 October 2026\\Version 0.2.0 --- unpublished working revision}
\begin{document}
\maketitle
\begin{abstract}
We observe only the isomorphism class of the strict coordinatewise order
of $n$ independent points in the unit square. Their unknown density is
smooth, lies between $1/2$ and $3/2$, has uniform marginals, and is
Euclidean Lipschitz with constant at most two. We construct one measurable
order-only estimator whose full-square supremum error, modulo one global
coordinate exchange, is at most
$\min\{1/2,650(\log n/n)^{1/4}\}$ with probability at least $19/20$,
uniformly over this class and for every $n\ge2$.
The principal reduction uses observed predecessor and successor counts:
every retained point has reliable reconstructed coordinates, while every
point in a sufficiently deep cell is retained. This removes a global
discarded-mass penalty from local density estimation. We also obtain a
finite inverse bound for the actual order laws. Smooth marginal-preserving
Gaussian alternatives give both a finite two-point obstruction and a
matching logarithmic lower bound: for $n\ge2^{64}$, every eligible
estimator, including independent randomization, has strict error greater
than $(\log n/n)^{1/4}/8192$ with probability at least $1/2$ under some
density in the same class. Thus the uniform $95\%$ minimax radius has
rate $(\log n/n)^{1/4}$ up to constants. The upper, law, lower and
minimax-radius endpoints are Lean-certified. The lower argument already
holds with coordinate data. Efficient computation, sharp constants and
useful practical calibration remain open.
\end{abstract}

\section{Experiment, results and contribution}
\label{sec:results}
Let $D=[0,1]^2$. The class $\K$ consists of functions smooth on an open
neighborhood of $D$ satisfying, for $x,y\in D$ and $u,v\in[0,1]$,
\begin{align}
 \tfrac12\le\rho(x)\le\tfrac32,\qquad
 |\rho(x)-\rho(y)|&\le2|x-y|_2, \label{eq:class}\\
 \int_0^1\rho(u,t)\dd t&=1,\qquad
 \int_0^1\rho(t,v)\dd t=1. \nonumber
\end{align}
The norm in the Lipschitz condition is Euclidean. In particular
$Q_\rho(\dd u\dd v)=\rho(u,v)\dd u\dd v$ on $D$ is a probability measure.
Every density is smooth, but no common bound on higher derivatives is
assumed. Uniform marginals fix the coordinate gauge.

From $Z_i=(U_i,V_i)$, $1\le i\le n$, drawn independently from $Q_\rho$,
we observe the directed relation
\[
 i\prec j\quad\Longleftrightarrow\quad U_i<U_j\ \hbox{and}\ V_i<V_j
\]
only up to vertex relabeling. Write $\mathcal O_n$ for this finite
observation space and $\pi_n^\rho$ for its actual pushforward law.
One may include all directed Boolean relation codes in $\mathcal O_n$
and give invalid codes a fixed fallback; this agrees with the formal
implementation. No coordinate order, rank, embedding or extra sample
is observed. The direction is retained: time reversal is not quotiented.

Let $T(u,v)=(v,u)$ and $G=\{\mathrm{id},T\}$. For bounded functions put
\[
 L(b,\rho)=\min_{S\in G}\norm{b-\rho\circ S},\qquad
 \dc(\rho,\sigma)=L(\rho,\sigma).
\]
All norms are over the full closed square. One $S$ must work at every
point. More generally, success at radius $a$ means
\begin{equation}\label{eq:success}
 \exists S\in G\ \ \forall x\in D,\quad |b(x)-\rho(Sx)|\le a.
\end{equation}
For unbounded outputs we use this definition directly (or an extended
supremum), avoiding any boundedness assumption in the lower theorem.
Total variation is $\TV(P,Q)=\sup_A|P(A)-Q(A)|$.

\begin{theorem}[One-order upper guarantee]\label{thm:upper}
For every integer $n\ge2$ there exists a measurable estimator
$\widehat\rho_n:\mathcal O_n\to\{b:D\to[1/2,3/2]\}$, independent of
$\rho$, whose outputs are measurable functions, such that for every
$\rho\in\K$,
\begin{equation}\label{eq:upper}
 \pi_n^\rho\{L(\widehat\rho_n,\rho)\le R_n\}\ge\frac{19}{20},
 \qquad R_n=\min\left\{\frac12,650\left(\frac{\log n}{n}\right)^{1/4}\right\}.
\end{equation}
\end{theorem}
The finite observation space makes the estimator and its deterministic
success events measurable. The construction below also establishes
measurability in the query point. Outputs need not be smooth, integrate
to one or have uniform marginals; \emph{density estimator} refers to the
target being estimated, not membership of the output in $\K$.

\begin{corollary}[Actual-law inverse]\label{cor:law}
For $\rho,\sigma\in\K$ and $n\ge2$,
\begin{equation}\label{eq:law}
 \dc(\rho,\sigma)\le
 \min\left\{1,\,
 1300\left(\frac{\log n}{n}\right)^{1/4}
       +\frac{10}{9}\TV(\pi_n^\rho,\pi_n^\sigma)\right\}.
\end{equation}
\end{corollary}
This compares two probability laws; it is not an estimator of their
total variation from a single observed order.

\begin{theorem}[Finite randomized lower bound]\label{thm:lower}
For $n\ge1$ set
\[
 h=\min\{1/2,n^{-1/4}\},\quad f(t)=t e^{-t^2},\quad
 a_h(x)=f((x-1/2)/h),
\]
and $\rho_0=1$, $\rho_1(u,v)=1+2h a_h(u)a_h(v)$.
These two densities belong to $\K$. For any estimator from
$\mathcal O_n$, including an arbitrary independent probability-space
random seed, assume the two success events \eqref{eq:success} at radius
$h/6$ are measurable. Then
\begin{equation}\label{eq:lower}
 \max_{i=0,1}\Pr_{\rho_i}\{L(\widehat\rho,\rho_i)>h/6\}
 \ \ge\ \frac{1-B_n}{2}\ >\ \frac14,\qquad
 B_n=\frac12\sqrt{\exp(n\pi h^4/8)-1}.
\end{equation}
In particular, for $n\ge16$, every estimator with measurable success
events for all $\rho\in\K$ has, under some $\rho\in\K$, error strictly
greater than $n^{-1/4}/6$ with probability strictly greater than $1/4$.
\end{theorem}
For randomized estimators, measurability of evaluation at each point
alone need not imply measurability of an uncountable simultaneous
success event. The explicit success-event assumption is therefore part
of the theorem. Deterministic estimators on $\mathcal O_n$ satisfy it.

The preceding two-point result remains useful at smaller sample sizes.
The next result closes the logarithmic gap in the same model. Set
$a_n=(\log n/n)^{1/4}$.

\begin{theorem}[Logarithmic randomized lower bound]\label{thm:loglower}
For every integer $n\ge2^{64}$ and every estimator
$b:\mathcal O_n\times\Omega\to\R^D$, let $\xi$ be an arbitrary
probability law for its independent seed. Assume the success events
\eqref{eq:success} at radius $a_n/8192$ are measurable for every
$\rho\in\K$. Then some $\rho\in\K$ satisfies
\begin{equation}\label{eq:loglower}
 (\pi_n^\rho\otimes\xi)\{L(b,\rho)>a_n/8192\}\ge\frac12.
\end{equation}
The same conclusion holds at any smaller radius whose success events
are measurable. No regularity or class membership of the outputs is
required.
\end{theorem}

For a fixed seed law $\xi$, let $r_{95,\xi}(n)$ be the infimum of all
nonnegative radii $r$ for which one estimator has measurable success
events at $r$ and success probability at least $19/20$ for every
$\rho\in\K$. A one-point seed includes deterministic estimation.
\begin{corollary}[Fixed-confidence minimax rate]\label{cor:minimax}
For every probability seed space and $n\ge2^{64}$,
\begin{equation}\label{eq:minimax}
 \frac{a_n}{8192}\le r_{95,\xi}(n)
 \le\min\{1/2,650a_n\}.
\end{equation}
Consequently the same bounds hold if the infimum also ranges over
independent randomized procedures.
\end{corollary}
This is rate optimality up to constants at fixed confidence. It makes
no assertion about a sharp leading constant or expected risk. The
lower proof applies to full coordinate observations, while the upper
estimator uses order only. Both experiments therefore have the same
rate for this loss; equal finite-sample information or equal optimal
constants do not follow.

\subsection{Relation to earlier work}
The geometric input is the occupied-grid realizer theorem in the
author's reconstruction paper \cite{Reconstruction}. For completeness,
Section~\ref{sec:geometry} repeats its forcing argument and probability
calibration. Coordinate-pattern forcing is earlier work: Winkler's
1985 Theorem~7 and his dimension-two analysis
\cite{Winkler1985,Winkler} are direct precedents. In the independent
uniform-permutation model, Winkler's Lemmas~2.2--2.6 and Theorem~4.3
give stronger typical rank rigidity, up to a global axis exchange.
Nonconstant members of $\K$ do not give independent coordinate
permutations, so that flat-model conclusion is not imported here.

Copula density estimation from coordinate data or paired empirical
ranks is also established. Genest, Masiello and Tribouley
\cite{Genest} analyze wavelet estimators in integrated quadratic and
pointwise expected loss. Autin, Le Pennec and Tribouley \cite{Autin}
study adaptive thresholding and quadratic risk. These are distinct
observation and loss specifications from a finite-confidence,
order-only supremum bound.

The published Seck--Mamane result \cite{SeckMamane}, Theorem~1
(p.~40), concerns rank-based wavelets and almost-sure uniform
convergence under wavelet and Besov conditions. Its smoothness
condition is $(d+2)/2<t<N+1$; with
$2^{j_n}\asymp(n/\log n)^{1/(2t+d)}$, its displayed rate is
\[
 \left(\frac{2^{\,2(1+d)j_n}\log\log n}{n}\right)^{1/2}.
\]
Remark~3 (p.~43) distinguishes this from the optimal oracle rate.
The published theorem differs from the earlier 2023 preprint, so we
use the published statement. Swanepoel, Seck and Mamane \cite{Swanepoel}
obtain an almost-sure supremum rate via joint and marginal density
estimates and quantiles, with their own Besov, support and positive
marginal assumptions. Neither almost-sure statement is converted
here into a uniform finite $95\%$ radius. Smoothness of each $\rho$
does not supply uniform higher-regularity norms over $\K$.

Braun's qualitative reconstruction results \cite{Braun} concern
all-size chronology laws under stated geometric hypotheses, with
dimension at least three. The present fixed-gauge, two-dimensional
statistical experiment does not extend those theorems or settle
general spacetime reconstruction. The separate coordinate-gauge
paper \cite{Gauge} exhibits a higher-dimensional obstruction to a
proposed gauge; it is a different manuscript and result.

The candidate contribution here is the degree-filter/deep-cell
reduction in Section~\ref{sec:filter}, composed with geometric
recovery to obtain \eqref{eq:upper}, and its exact formal verification.
The new contribution in this revision is a matching lower bound in
the exact original class and quotient, and a certified minimax-radius
sandwich. Histogram bias, Chernoff bounds, likelihood products and
testing reductions are standard methods. In particular, the
second-moment many-hypothesis inequality used below is the $l=2$
case of the established power-divergence route in Guntuboyina
\cite{Guntuboyina}, Example~II.8, equation~(15), following the
finite-packing reduction in Section~I and Theorem~II.1. We write out
the elementary inequality and its seed and measurability details;
we do not claim a new general testing method. The bounded primary-source comparison
found no inspected theorem with the full experiment, class, quotient
and confidence specification of Theorem~\ref{thm:upper}; this is a
provisional priority assessment, not an exhaustive absence claim.

\section{Geometric recovery on a finite-probability event}
\label{sec:geometry}
Fix $m\ge16$, $r=1/m$. Let $E$ require every open grid cell of side $r$
to be occupied, and the empirical anchored distribution function
$C_n(s,t)=n^{-1}\#\{i:U_i\le s,V_i\le t\}$ to differ from its population
version $C(s,t)$ by at most $r$ at each grid vertex.
We also intersect with the probability-one event of distinct
coordinates in $(0,1)$ and no coordinate on any grid line.
The latter restrictions do not change any probability bounds.

\begin{lemma}[Grid probability and interpolation]\label{lem:grid}
Uniformly over $\rho\in\K$,
\begin{equation}\label{eq:geometricfailure}
 \Pr_\rho(E^c)\le m^2 e^{-nr^2/2}+2(m+1)^2e^{-2nr^2}.
\end{equation}
On $E$ the two marginal empirical errors at arbitrary thresholds in
$[0,1]$ are at most $2r$, the joint error is at most $3r$, and the
fraction of points outside $J=[4r,1-4r]^2$ is at most $20r$.
\end{lemma}
\begin{proof}
Every cell has probability at least $r^2/2$, so its empty probability
is at most $(1-r^2/2)^n\le e^{-nr^2/2}$.
At a fixed grid vertex, Hoeffding's inequality \cite{Hoeffding}
gives a two-tail probability at most $2e^{-2nr^2}$. A union bound
proves \eqref{eq:geometricfailure}. For clarity, the bounded-variable
MGF behind this step is
$\mathbb E e^{t(X-\mathbb EX)}\le e^{t^2/8}$ for $0\le X\le1$:
the second derivative of its log MGF is the variance under exponential
tilting, at most $1/4$, and the value and first derivative at zero
vanish. Independence, Markov's inequality and $t=4r$ yield one tail
$e^{-2nr^2}$; use $-t$ for the other.

At grid thresholds, each marginal is an anchored rectangle with the
other threshold $1$, so its error is at most $r$. Monotonicity and
uniform marginals add at most $r$ between adjacent grid thresholds.
For the joint CDF, enlarge or shrink the rectangle to a grid vertex.
The population difference is at most $2r$ because it lies in two
marginal strips. Adding the vertex error gives $3r$.
Each of the four marginal tails beyond $4r$ and $1-4r$ has empirical
mass at most $5r$. Their union contains all points outside $J$.
\end{proof}

A realizer is a pair of strict total orders whose intersection is
$\prec$. Its inclusive rank is
$\rank_L(i)=1+\#\{j:L(j,i)\}$. Distinct latent coordinates supply a
realizer, although neither latent total order is observed.

\begin{lemma}[Earlier realizer rigidity]\label{lem:rank}
On $E$, every realizer has one global exchange of its two orders
after which, for every $Z_i\in J$,
\begin{equation}\label{eq:rank}
 \left|\frac{\rank_{L_1}(i)}n-U_i\right|\le32r,\qquad
 \left|\frac{\rank_{L_2}(i)}n-V_i\right|\le32r.
\end{equation}
\end{lemma}
\begin{proof}
This is the finite theorem of \cite{Reconstruction}; we give its
argument to specify the dependency. For incomparable $x,y$ write
$Q(x,y)$ when $L_1(x,y)$ and $L_2(y,x)$. The relation $Q$ is
transitive. If incomparability edges $xy,xz$ share $x$ and $y,z$
are comparable, their orientations at $x$ agree; otherwise a
two-edge $Q$ path would make $y,z$ incomparable.

Occupancy provides
\begin{align*}
 A&\in(r,2r)\times(1-2r,1-r),\\
 B&\in(1-2r,1-r)\times(r,2r),\\
 W&\in(1-r,1)\times(1-3r,1-2r).
\end{align*}
For interior incomparable points $x,y$ with $u_x<u_y$ and
$v_x-v_y\ge3r$, choose an open grid row inside $(v_y,v_x)$.
Explicitly, $k=\lfloor v_y/r\rfloor+1$ has $kr>v_y$ and
$(k+1)r\le v_y+2r<v_x$; interior membership makes this row admissible.
The first cell in that row supplies $z$ with $0<u_z<r$.
The incomparability edges
\[
 xy,\quad zy,\quad Ay,\quad AW,\quad AB
\]
have comparable opposite endpoints at each step:
$z\prec x$, $z\prec A$, $y\prec W$, $B\prec W$.
Thus $Q(x,y)$ has the same truth value as $Q(A,B)$.
Exchanging the two orders once makes $Q(A,B)$ agree with increasing
latent $u$. All such separated interior incomparable pairs then
agree with the latent realizer; comparable pairs already agree.

For a fixed interior point $Z_i$, a rank disagreement can therefore
involve only a point outside $J$ or one in $|V_j-V_i|<3r$.
The latter strip has empirical mass at most
$6r+4r=10r$ by Lemma~\ref{lem:grid}; both endpoints lie in $[0,1]$.
There are at most $30rn$ possible disagreements, bounding each
rank difference by $30rn$. Finally, the inclusive true rank divided
by $n$ equals the marginal empirical CDF at the point and is within
$2r$ of its coordinate. This proves \eqref{eq:rank}.
\end{proof}
The orientation can depend on the sample and realizer. It is a single
choice for all retained points and all subsequent queries.

\section{Observed degrees and the deep-cell reduction}
\label{sec:filter}
Let $h=\sqrt r$ and retain precisely
\begin{equation}\label{eq:filter}
 I=\{i:d_i^->6rn,\ d_i^+>6rn\},
\end{equation}
where $d_i^-=\#\{j:j\prec i\}$ and $d_i^+=\#\{j:i\prec j\}$.
These counts use only the observed relation.

\begin{lemma}[Two inclusions]\label{lem:deep}
If $E$ holds and $1/n\le r$, then
\[
 \{i:Z_i\in[7h,1-7h]^2\}\ \subseteq\ I\
 \subseteq\ \{i:Z_i\in J\}.
\]
\end{lemma}
\begin{proof}
If $U_i<4r$ or $V_i<4r$, all predecessors are in the corresponding
empirical lower marginal tail, of mass at most $5r$.
If $U_i>1-4r$ or $V_i>1-4r$, all successors lie in an upper tail
of mass at most $5r$. Such points fail \eqref{eq:filter}.

For the converse use the exact strict-count identities
\begin{align}
 d_i^-/n&=C_n(U_i,V_i)-1/n,\label{eq:self}\\
 d_i^+/n&=1-F_{U,n}(U_i)-F_{V,n}(V_i)+C_n(U_i,V_i).
 \label{eq:successor}
\end{align}
Coordinate injectivity makes the self-subtraction in \eqref{eq:self}
exact; inclusion-exclusion gives \eqref{eq:successor}.
At a deep point, southwest and northeast population rectangle
masses are each at least $(7h)^2/2=49r/2$, by the lower density bound.
The first empirical rectangle differs by at most $3r$, and the
second by at most $2r+2r+3r=7r$. Consequently
\[
 d_i^-/n\ge49r/2-3r-1/n\ge41r/2>6r,\quad
 d_i^+/n\ge49r/2-7r=35r/2>6r.
\]
No independence of degrees or retained points is used.
\end{proof}

Now suppose $m\ge65536$ and define
\begin{equation}\label{eq:parameters}
 k=\lfloor\sqrt m\rfloor,\quad H=8h,\quad s=32r,\quad
 w=\frac{1-2H}{k}.
\end{equation}
Partition the inner square into $k^2$ cells of side $w$, using
intervals $(a,a+w]$ for counts. For a fixed realizer let
$q_i=(\rank_{L_1}(i)/n,\rank_{L_2}(i)/n)$ and set
\begin{equation}\label{eq:histogram}
 b_R=\frac{\#\{i\in I:q_i\in R\}}{nw^2}.
\end{equation}
The denominator is the original $n$, not the number retained.
For $R=(a,a+w]\times(c,c+w]$, write
\[
 R^-=(a+s,a+w-s]\times(c+s,c+w-s],\qquad
 R^+=(a-s,a+w+s]\times(c-s,c+w+s].
\]
Fixed boundary lines are null for latent points. These half-open
conventions also make the deterministic inclusions below valid at
their relevant endpoints.

\begin{lemma}[Mesh and cell sandwich]\label{lem:sandwich}
The parameters \eqref{eq:parameters} satisfy
\begin{equation}\label{eq:mesh}
 0<h\le\frac1{256},\quad \frac{15}{16}h\le w\le2h,\quad
 2s\le\frac h4<w,\quad H-s\ge\frac{63}8h>7h.
\end{equation}
If $E$ holds and $1/n\le r$, one $S\in G$ simultaneously gives,
for every cell,
\begin{equation}\label{eq:sandwich}
 \frac{N(SR^-)}n\le
 \frac{\#\{i\in I:q_i\in R\}}n\le
 \frac{N(SR^+)}n,
\end{equation}
where $N(A)=\#\{i:Z_i\in A\}$.
\end{lemma}
\begin{proof}
We have $\sqrt m/2\le k\le\sqrt m$, $k^2\le m$ and
$15/16\le1-16h<1$, giving the width inequalities.
The other bounds follow from $s=32h^2$ and $h\le1/256$.
Every expanded cell, including its closure, lies in
$[7h,1-7h]^2$ by the final margin bound.

Fix the one alignment in Lemma~\ref{lem:rank}; equivalently
$|q_i-SZ_i|_\infty\le s$ for every retained point.
If $SZ_i\in R^-$, Lemma~\ref{lem:deep} retains $i$, and its coordinate
error places $q_i$ in $R$. If $q_i\in R$ and $i$ is retained, the same
error puts $SZ_i$ in $R^+$. These implications prove
\eqref{eq:sandwich}. Both squares and the degree filter are invariant
under transpose, so the same $S$ works throughout.
\end{proof}
This is the key reduction: local true points are all retained.
There is no global discarded-count term divided by the cell area.

\section{Local concentration and full-square estimation}
\label{sec:upperproof}
\begin{lemma}[Small-mass concentration]\label{lem:concentration}
For independent Bernoulli variables of common mean $p\le8h^2$,
where $0<h\le4$, their average $\overline X$ satisfies
\begin{equation}\label{eq:chernoff}
 \Pr\{|\overline X-p|>4h^3\}\le2e^{-nh^4/2}.
\end{equation}
\end{lemma}
\begin{proof}
For $|t|\le1$, the exponential series gives
$e^t-1-t\le t^2\sum_{j\ge2}1/j!\le t^2$.
Here $e<3$, as also shown in Section~\ref{sec:lowerproof}.
Thus
\[
 \mathbb E e^{t(X-p)}
 =e^{-tp}\{1+p(e^t-1)\}
 \le e^{p(e^t-1-t)}\le e^{pt^2}.
\]
Multiply these MGFs, apply Markov's inequality and choose $t=h/4$.
The exponent for tolerance $4h^3$ is at most
$-n(h/4)(4h^3)+n(8h^2)(h/4)^2=-nh^4/2$.
The same calculation with $-t$ proves the other tail.
\end{proof}

The family of rectangles $SR^\pm$, over every cell and both $S\in G$,
is fixed before sampling. Each member satisfies
\begin{equation}\label{eq:mass}
 Q_\rho(SR^\pm)\le\tfrac32(w+2s)^2
 \le\tfrac32(9h/4)^2=\tfrac{243}{32}h^2<8h^2.
\end{equation}
Counting the two signs and two orientations gives at most $4k^2$
rectangles. Let $F$ be the event that all their empirical masses
are within $\varepsilon=4h^3$ of their true masses. Then
\begin{equation}\label{eq:localfailure}
 \Pr(F^c)\le8k^2e^{-nh^4/2}.
\end{equation}
Some rectangles duplicate one another under transpose; counting both
orientations is the conservative bound used by the formal proof.
We do not assume that recovered points are independent, or that $E$
and $F$ are independent.

\begin{lemma}[Cell error and boundary extension]\label{lem:bias}
With $m\ge65536$ and $1/n\le r$, on $E\cap F$, clip each $b_R$ to
$[1/2,3/2]$. Clamp each query coordinate into $[H,1-H]$ and return
the clipped value of a cell whose closure contains the clamped point.
This gives, under the same $S$, a full-square error at most $270h$.
\end{lemma}
\begin{proof}
Put $\rho^S=\rho\circ S$ and $p_R=\int_R\rho^S$.
The expansion shell has area $4ws+4s^2$, so its mass is at most
$6ws+6s^2$. The contraction loss is no larger.
The count sandwich and event $F$ give
\begin{equation}\label{eq:cell}
 |b_R-p_R/w^2|\le 6s/w+6(s/w)^2+4h^3/w^2.
\end{equation}
At every point $y$ in the closed cell, the difference between
$\rho^S(y)$ and its cell average is at most $2w$.
Indeed, for a uniform point $X$ in the cell,
$\mathbb E|X-y|_2^2\le2w^2/3$, with the largest value at a corner;
Cauchy--Schwarz and the Euclidean Lipschitz constant give
$2\mathbb E|X-y|_2\le2w$.

Clipping cannot increase error relative to a number in $[1/2,3/2]$.
Clamping moves a query point by at most $\sqrt2 H$, so its density
cost is at most $2\sqrt2H<3H=24h$. Using \eqref{eq:mesh} in
\eqref{eq:cell}, we obtain respectively
\[
 6s/w\le\frac{3072}{15}h,\quad
 6(s/w)^2\le\frac{6144}{225}h,\quad
 4h^3/w^2\le\frac{1024}{225}h .
\]
The sum of these three constants is $53248/225$.
Adding $2w\le4h$ and $24h$ gives
$(59548/225)h<270h$, uniformly on the closed square.
\end{proof}

To fix the query convention exactly, set
\[
 j(x)=\min\left\{\left\lfloor\frac{\operatorname{clamp}_{[H,1-H]}(x)-H}{w}
                    \right\rfloor,k-1\right\}.
\]
Use cell $(j(u),j(v))$. Internal query edges select the cell to their
right, while count cells are open on the left and closed on the right.
Lemma~\ref{lem:bias} holds on each cell closure, so this distinction
does not affect any query, including corners. The coordinatewise
convention commutes with transpose. Finite cell selection and clamping
also prove query measurability.

\subsection{The selector and every sample-size branch}
For $n\ge2$, let
\begin{equation}\label{eq:outermesh}
 x_n=\sqrt{\frac{n}{8\log n}},\qquad m=\lfloor x_n\rfloor .
\end{equation}
If $m<65536$, return the flat function $1$. Otherwise put $h=m^{-1/2}$.
If $270h\ge1/2$, again return $1$. These decisions depend only on $n$.
In the remaining branch use the construction above.

Choose once, on the finite space $\mathcal O_n$, one labeled
representative of each isomorphism class and one realizer whenever
one exists. Finite enumeration is one possible specification; the Lean
selector uses finite classical choices. Return $1$ if there is no
realizer. On actual samples there is a realizer almost surely.
Relabeling permutes the predecessor and successor sets and the rank
counts, so the construction transports to the selected representative.
The proofs hold for every realizer, with its own single global
alignment. Thus the choice uses no unknown $\rho$, and the guarantee
descends to the actual unlabeled observation. No efficient algorithm
for making these choices is asserted.

\begin{proof}[Proof of Theorem~\ref{thm:upper}]
In the active branch $m\ge65536$, and
\[
 nr^2=n/m^2\ge8\log n,\quad
 m^2\le n/8,\quad n\ge4m^2,\quad 1/n\le r.
\]
Here $n\ge4m^2$ follows first from $\log n\ge\log2>1/2$;
this implies $n\ge4$ in this branch, hence $\log n>1$ and $m^2\le n/8$.
Since $(m+1)^2\le4m^2$, Lemma~\ref{lem:grid} gives
\begin{equation}\label{eq:finalgeometric}
 \Pr(E^c)\le\frac1{8n^3}+\frac1{n^{15}}.
\end{equation}
Also $h^4=r^2$, $k^2\le m\le n$, so \eqref{eq:localfailure} gives
$\Pr(F^c)\le8/n^3$. Already for $n\ge8$,
\[
 \frac1{8n^3}+\frac1{n^{15}}+\frac8{n^3}<\frac1{20}.
\]
Lemma~\ref{lem:bias} consequently supplies the $270h$ radius with
probability at least $19/20$. This is the certified intermediate
Chernoff budget, not the sharper Bernstein budget of an earlier
working derivation.

For $m\ge65536$, $\lfloor x_n\rfloor\ge x_n/2$, whence
\[
 270h\le270\sqrt2\,8^{1/4}
        \left(\frac{\log n}{n}\right)^{1/4}
 <650\left(\frac{\log n}{n}\right)^{1/4}.
\]
The numerical comparison is $32\cdot270^4<650^4$.
In the active branch $270h<1/2$ as well, giving the minimum in
\eqref{eq:upper}.

If $m\ge65536$ but $270h\ge1/2$, the same scale comparison gives
$R_n=1/2$. If $m<65536$, then $x_n<65536$ and
\[
 650\left(\frac{\log n}{n}\right)^{1/4}
 >\frac{650}{8^{1/4}\,256}>\frac12,
\]
so again $R_n=1/2$. In both fallback branches the flat function has
error at most $1/2$ deterministically, by \eqref{eq:class}.
This covers every $n\ge2$ and completes the selector proof.
\end{proof}

\begin{proof}[Proof of Corollary~\ref{cor:law}]
Let $A_\rho,A_\sigma$ be the two success sets of the same estimator
at radius $R_n$, and put $\delta=\TV(\pi_n^\rho,\pi_n^\sigma)$.
Under $\pi_n^\sigma$, their intersection has probability at least
$19/20-\delta+19/20-1=9/10-\delta$.
If $\delta<9/10$, this intersection is nonempty. At any observation
in it, the two global alignments and the triangle inequality imply
$\dc(\rho,\sigma)\le2R_n\le1300(\log n/n)^{1/4}$.
If $\delta\ge9/10$, the class bounds give
$\dc(\rho,\sigma)\le1\le(10/9)\delta$.
In all cases $\dc\le1$, yielding \eqref{eq:law}.
\end{proof}

\section{Smooth alternatives and testing}
\label{sec:lowerproof}
\begin{lemma}[Original-class membership and separation]\label{lem:alternatives}
For $0<h\le1/2$, the two alternatives in Theorem~\ref{thm:lower}
belong to $\K$, are invariant under transpose, and have
$\dc(\rho_0,\rho_1)=h/e>h/3$.
\end{lemma}
\begin{proof}
The profile $f(t)=te^{-t^2}$ is smooth and odd, with
\[
 \sup|f|=\frac1{\sqrt{2e}}<\frac12,\qquad
 f'(t)=(1-2t^2)e^{-t^2},\qquad \sup|f'|\le1.
\]
The first maximum occurs at $|t|=1/\sqrt2$. For the derivative,
on $t^2\le1/2$ the bound is immediate; on $t^2\ge1/2$ the maximum
of $(2t^2-1)e^{-t^2}$ is $2e^{-3/2}<1$, at $t^2=3/2$.
Since $a_h(1-x)=-a_h(x)$, its integral on $[0,1]$ vanishes exactly.
Consequently both marginal integrals of $\rho_1$ equal one.
Also $|\rho_1-1|\le h/2\le1/4$, and each partial derivative has
magnitude at most $2\sup|f'|\sup|f|\le1$.
The Euclidean gradient bound $\sqrt2<2$, integrated along segments
in $D$, proves \eqref{eq:class}. Smoothness holds on all of $\R^2$.

Both densities are transpose invariant. At
$u=v=1/2+h/\sqrt2$, which lies in $D$, the deviation is $h/e$.
The profile maximum bounds the deviation everywhere by the same
number, proving exact quotient separation.
Finally $e<3$: for $j\ge2$, $j!\ge2^{j-1}$, with strict inequality
for $j\ge3$, so $\sum_{j\ge2}1/j!<1$. Thus $h/e>h/3$.
\end{proof}

\begin{lemma}[Likelihood product and order data processing]\label{lem:divergence}
For the alternatives above and every $n\ge1$,
\[
 \TV(\pi_n^{\rho_0},\pi_n^{\rho_1})
 \le\frac12\sqrt{\exp(n\pi h^4/8)-1}.
\]
\end{lemma}
\begin{proof}
With $Q_0$ uniform on $D$,
\[
 c:=\chi^2(Q_1,Q_0)
 =\int_D(\rho_1-1)^2
 =4h^2\left(\int_0^1 a_h(x)^2\dd x\right)^2
 \le\frac{\pi}{8}h^4.
\]
To check the constant, substitute $t=(x-1/2)/h$ and enlarge the
nonnegative integral to the real line:
\[
 \int_{\R}t^2e^{-2t^2}\dd t=\frac14\sqrt{\frac\pi2},
 \qquad
 \left(\int_{\R}t^2e^{-2t^2}\dd t\right)^2=\frac\pi{32}.
\]
Indeed, squaring the Gaussian integral and using polar coordinates
gives $\int_\R e^{-2t^2}\dd t=\sqrt{\pi/2}$.
Integrating the derivative of $te^{-2t^2}$ on expanding symmetric
intervals, whose boundary term tends to zero, gives the factor $1/4$.

The actual likelihood ratio of $Q_1^{\otimes n}$ against
$Q_0^{\otimes n}$ is $W(z_1,\ldots,z_n)=\prod_i\rho_1(z_i)$.
Fubini, $\int\rho_1=1$ and $\int\rho_1^2=1+c$ give
\[
 \int W=1,\qquad \int W^2=(1+c)^n,\qquad
 \int(W-1)^2=(1+c)^n-1\le e^{n\pi h^4/8}-1.
\]
Therefore Cauchy--Schwarz gives
$\TV(Q_0^{\otimes n},Q_1^{\otimes n})
=\frac12\int|W-1|\le\frac12\sqrt{\int(W-1)^2}$.
The sampled directed relation is a finite collection of strict Borel
inequalities; passing to an isomorphism class is a finite map.
Pulling back an order event proves that its probability difference
cannot exceed coordinate-data TV. This proves the claim for the
actual unlabeled order laws.
\end{proof}

\begin{proof}[Proof of Theorem~\ref{thm:lower}]
Set $r_*=h/6$. By Lemma~\ref{lem:alternatives}, the closed success
sets at radius $r_*$ for $\rho_0,\rho_1$ are disjoint:
otherwise the triangle inequality after composing their global
alignments would give $\dc(\rho_0,\rho_1)\le2r_*=h/3$.
This argument applies to arbitrary outputs; it uses only the two
all-point success inequalities.

Write $P_i=\pi_n^{\rho_i}$. If a seed has probability law $\xi$,
independent of the observation under either hypothesis, every
measurable event $A\subseteq\mathcal O_n\times\Omega$ has
\[
 |(P_0\otimes\xi)(A)-(P_1\otimes\xi)(A)|
 \le\int |P_0(A_\omega)-P_1(A_\omega)|\dd\xi(\omega)
 \le\TV(P_0,P_1).
\]
Sections are measurable and the sum over the finite observation
space justifies the integral without additional output regularity.
Decide hypothesis zero exactly on its success event $A$.
Under zero the test error is strict estimation failure; under one,
disjointness makes test error a subset of strict estimation failure.
The sum of these two risks is at least
\[
 (P_0\otimes\xi)(A^c)+(P_1\otimes\xi)(A)
 =1-\{(P_0\otimes\xi)(A)-(P_1\otimes\xi)(A)\}
 \ge1-B_n.
\]
Their maximum is therefore at least $(1-B_n)/2$.

For the selected $h$, $nh^4\le1$: if $n\le16$ use $h=1/2$;
if $n\ge16$ use $h=n^{-1/4}$. Since $\pi<4$ and
$e^{1/2}<2$, we have $B_n<1/2$. The latter exponential inequality
also follows from $e<3<4$. Thus the finite lower bound is strictly
greater than $1/4$. Both alternatives are in $\K$, and for $n\ge16$
the radius is exactly $n^{-1/4}/6$, proving the final assertion.
\end{proof}

\subsection{What follows for confidence bands}
Suppose measurable simultaneous bands $[\ell(x),u(x)]$ have
$\ell\le u$, with coverage meaning containment of $\rho\circ S$ at
every $x\in D$ for one $S\in G$. Assume the relevant band and midpoint
success events are measurable. If the maximum width
$W=\sup_D(u-\ell)$ is at most $2r_*$ on every output, coverage implies
midpoint loss at most $r_*$. Such a band cannot have uniform $95\%$
coverage at this width. If the width is random, the valid implication is
\begin{equation}\label{eq:band}
 \max_{i=0,1}\left[
 \Pr_{\rho_i}\{\text{coverage fails}\}
 +\Pr_{\rho_i}\{W>2r_*\}\right]\ge\frac{1-B_n}{2}.
\end{equation}
Indeed midpoint failure is contained in the union of these two
events, and Theorem~\ref{thm:lower} applies to the midpoint.
This does not exclude occasional narrow reports or assert anything
about coverage conditional on accepting a report. Equation~\eqref{eq:band}
is a written consequence, not a separately exported Lean band theorem.

\section{A matching logarithmic obstruction}
\label{sec:logproof}
We retain $f(t)=te^{-t^2}$ from Section~\ref{sec:lowerproof}, but
translate and center it to allow a growing family of alternatives.

\begin{lemma}[Centered profiles]\label{lem:logprofile}
For $h>0$ and $c\in\R$, define
\[
 g_{h,c}(x)=f((x-c)/h),\quad
 \mu_{h,c}=\int_0^1 g_{h,c}(x)\dd x,\quad
 p_{h,c}(x)=g_{h,c}(x)-\mu_{h,c}.
\]
The profile is smooth on $\R$, has zero integral on $[0,1]$, and
\begin{align}
 |\mu_{h,c}|&\le h/2,& |p_{h,c}(x)|&\le1/2+h/2,\label{eq:logprofile}\\
 |p_{h,c}(x)-p_{h,c}(y)|&\le|x-y|/h,&
 \int_0^1p_{h,c}^2&\le h.\nonumber
\end{align}
If $h\le1/16$, then $p_{h,c}(c+h/2)\ge1/4$ and
$|p_{h,c}(x)|\le1/8$ whenever $|(x-c)/h|\ge4$.
\end{lemma}
\begin{proof}
An antiderivative of $g_{h,c}$ is
$-(h/2)e^{-((x-c)/h)^2}$, so
\[
 \mu_{h,c}=\frac h2\bigl(e^{-(c/h)^2}-e^{-((1-c)/h)^2}\bigr).
\]
Both exponentials lie in $[0,1]$, proving the mean bound.
The bounds $|f|\le1/2$ and $|f'|\le1$ proved earlier give the
absolute and Lipschitz bounds. Centering gives exact zero mean and
\[
 \int_0^1p_{h,c}^2
 =\int_0^1g_{h,c}^2-\mu_{h,c}^2
 \le h\int_\R f(t)^2\dd t
 =\frac{h\sqrt\pi}{4\sqrt2}\le h.
\]
All interval integrands are continuous; the whole-line integral is
the finite Gaussian moment evaluated in Lemma~\ref{lem:divergence}.
Since $e^{-1/4}\ge3/4$, $f(1/2)\ge3/8$; subtracting a mean of
absolute value at most $h/2$ proves the peak bound.
For $|t|\ge4$,
\[
 e^{t^2}\ge(1+t^2/2)^2\ge t^4/4\ge16|t|,
\]
using $e^s\ge1+s$ and $|t|^3\ge64$.
Thus $|f(t)|\le1/16$, and centering gives
$|p_{h,c}(x)|\le1/16+h/2\le3/32\le1/8$ at the stated tail points.
\end{proof}

\begin{lemma}[An admissible finite packing]\label{lem:logpacking}
For an integer $m\ge1$, put
\begin{equation}\label{eq:logfamily}
 h=\frac1{16m},\quad c_j=\frac14+\frac{j+1/2}{2m},\quad
 \rho_j(u,v)=1+\frac h2 p_{h,c_j}(u)p_{h,c_j}(v),
 \qquad 0\le j<m.
\end{equation}
All $\rho_j$ belong to $\K$ and are invariant under transpose.
For $j\ne k$ they have quotient separation at least $3h/128$.
Consequently their success sets at any radius $r\le h/256$ are
pairwise disjoint, even for arbitrary estimator outputs.
\end{lemma}
\begin{proof}
We have $0<h\le1/16$ and $|p_{h,c_j}|\le1$. Smoothness holds
globally. Zero profile mean makes each marginal integral exactly one,
and $|\rho_j-1|\le h/2\le1/32$ implies the required range.
Expanding a product difference and applying Lemma~\ref{lem:logprofile}
gives, for $z,w\in D$,
\[
 |\rho_j(z)-\rho_j(w)|
 \le\frac{|z_1-w_1|+|z_2-w_2|}{2}\le2|z-w|_2.
\]
This verifies the original Euclidean Lipschitz condition, without
a uniform higher-derivative assumption.

The centers lie in $[1/4,3/4]$ and satisfy $|c_j-c_k|\ge8h$.
The point $x_j=c_j+h/2$ lies in $[0,1]$. At $x_j$ the own profile
is at least $1/4$, whereas every other profile has absolute value
at most $1/8$, since $|(x_j-c_k)/h|\ge8-1/2\ge4$.
At the diagonal witness $z_j=(x_j,x_j)$ it follows that
\begin{equation}\label{eq:logseparation}
 \rho_j(z_j)-\rho_k(z_j)
 \ge\frac h2\left(\frac1{16}-\frac1{64}\right)=\frac{3h}{128}.
\end{equation}
Transpose fixes each alternative, so taking the quotient does not
decrease their supremum separation. Two successes would, at $z_j$,
give a difference at most $2r\le h/128$, contradicting
\eqref{eq:logseparation}. This proof uses only pointwise inequalities.
\end{proof}

\begin{lemma}[Product and independent-seed moments]\label{lem:logmoment}
Let $Q_0$ be uniform measure on $D$, and $Q=Q_0^{\otimes n}\otimes\xi$.
The likelihood of $Q_{\rho_j}^{\otimes n}\otimes\xi$ against $Q$ is
$L_j(z,\omega)=\prod_{i=1}^n\rho_j(z_i)$, and
\begin{equation}\label{eq:logmoment}
 \int L_j\dd Q=1,\qquad
 \int L_j^2\dd Q\le B:=e^{nh^4/4}.
\end{equation}
\end{lemma}
\begin{proof}
Zero marginal means cancel the linear term in $\rho_j^2$.
Fubini and Lemma~\ref{lem:logprofile} give
\[
 \int_D(\rho_j-1)^2
 =\frac{h^2}{4}\left(\int_0^1p_{h,c_j}^2\right)^2
 \le\frac{h^4}{4},\qquad \int_D\rho_j^2\le1+h^4/4.
\]
The densities are positive and bounded on $D$, so product likelihoods
and their squares are integrable. Independence gives
$\int L_j^2\le(1+h^4/4)^n\le e^{nh^4/4}$ and $\int L_j=1$.
Integration over an independent probability seed contributes the
factor one. In particular, for every measurable event $A$ the actual
sample-times-seed probability is $\int_A L_j\dd Q$.
\end{proof}

\begin{lemma}[Disjoint-event testing]\label{lem:manytesting}
Let $Q$ be a probability measure and $P_0,\ldots,P_{m-1}$ probability
measures with integrable likelihoods $L_j$ and $\int L_j^2\dd Q\le B$.
For pairwise disjoint measurable sets $A_j$,
\begin{equation}\label{eq:manytesting}
 \sum_{j=0}^{m-1}P_j(A_j)\le\sqrt{mB}.
\end{equation}
If $m\ge1$ and $B\le m/4$, some $P_j(A_j^c)\ge1/2$.
\end{lemma}
\begin{proof}
Set $F=\sum_j1_{A_j}L_j$. Disjointness implies
$F^2=\sum_j1_{A_j}L_j^2$ pointwise. Each summand is integrable,
and Cauchy--Schwarz on the probability space gives
\[
 \sum_jP_j(A_j)=\int F\dd Q
 \le\int|F|\dd Q\le\sqrt{\int F^2\dd Q}
 \le\sqrt{\sum_j\int L_j^2\dd Q}\le\sqrt{mB}.
\]
If $B\le m/4$, average success is at most $1/2$. Taking complements
under the probability measures proves the final assertion.
\end{proof}

\begin{lemma}[Explicit integer scale]\label{lem:logscale}
For $n\ge2^{64}$ put $t=\log n$, $x=(n/t)^{1/4}$, $m=\lceil x\rceil$
and $h=1/(16m)$. Then $m\ge16$, $x\le m\le2x$,
$e^{nh^4/4}\le m/4$, and $a_n/8192\le h/256$.
\end{lemma}
\begin{proof}
The inequality $\log2\ge1/2$ gives $t\ge32$.
For $t\ge0$, $e^{t/2}\ge(1+t/4)^2\ge t$, where the last
difference equals $(t-4)^2/16$. Taking logarithms yields
$\log t\le t/2$, hence
$\log x=(t-\log t)/4\ge t/8\ge4$.
Since $e\ge2$, $x\ge e^4\ge16$. The ceiling therefore satisfies
$x\le m<x+1\le2x$, and $n/m^4\le t$. Thus
\[
 \frac{nh^4}{4}=\frac{n/m^4}{262144}
 \le\frac t{16}\le\frac{\log m}{2}.
\]
Exponentiating gives $B\le\sqrt m\le m/4$ since $m\ge16$.
Finally $a_n=1/x$ and
$h/256=1/(4096m)\ge1/(8192x)=a_n/8192$.
\end{proof}

\begin{proof}[Proof of Theorem~\ref{thm:loglower}]
Choose the family and scale of Lemmas~\ref{lem:logpacking}
and~\ref{lem:logscale}. Pull back each success event at the target
radius through the map from coordinates and seed to unlabeled order
and seed. The map is measurable, since it consists of finitely many
strict coordinate comparisons and a finite quotient map, with the
identity on the seed. The pullbacks $A_j$ are therefore measurable
and pairwise disjoint. Their actual probabilities are the likelihood
integrals in Lemma~\ref{lem:logmoment}.
Lemma~\ref{lem:manytesting} and $B\le m/4$ give an alternative whose
complementary event has probability at least $1/2$.
The complement is strict failure of \eqref{eq:success}, including for
unbounded outputs. Applying the argument directly at the requested
radius uses no measurability assumption at an auxiliary radius.
The same argument works at any smaller measurable radius.
It also works for a coordinate-data estimator by using its success
sets directly, without the order-map pullback.
\end{proof}

\begin{proof}[Proof of Corollary~\ref{cor:minimax}]
For a fixed seed space, Theorem~\ref{thm:upper} supplies an estimator
that ignores its seed, so its success probability and measurability
are unchanged. The set defining $r_{95,\xi}(n)$ is nonempty and
bounded below by zero. Any achievable radius at most $a_n/8192$
would contradict Theorem~\ref{thm:loglower}: failure would be at least
$1/2$ for some density but at most $1/20$ for all densities.
Thus every achievable radius exceeds $a_n/8192$, and the infimum
is at least this number. The upper bound follows from the achievable
radius $\min\{1/2,650a_n\}$. Both constants are independent of the
seed law, which also permits the infimum over randomized procedures.
\end{proof}

\section{Formal verification and reproducibility}
\label{sec:formal}
All main endpoints have exact-source remote GCP acceptance.
Declarations below live in namespace \path{QuantyraNullCone}; a prefix
\path{InDensityClass} indicates its namespace within that root.
The companion \path{formal-map.json} supplies full identifiers, source
paths and the retained exact-type/axiom evidence.

\begin{description}
\item[Upper theorem.]
\path{fourth_root_order_only_estimator} and
\path{InDensityClass.fourth_root_density_estimation} in
\path{DegreeEstimator.lean} prove Theorem~\ref{thm:upper} and its
concrete selector instance. The estimator is quantified before the
unknown density, and no good-event hypothesis remains in the endpoint.
\item[Geometric and local reduction.]
\path{Cumulative.lean} supplies earlier coordinate rigidity.
\path{DegreeFilter.lean}, \path{DegreeMass.lean} and
\path{DegreeCells.lean} prove the strict-count identities,
deep retention and one-orientation count sandwich.
\path{BernoulliMGF.lean}, \path{SmallMassConcentration.lean} and
\path{DegreeProbability.lean} supply the actual iid concentration.
\item[Histogram and all-$n$ assembly.]
\path{DegreeCellBias.lean}, \path{InnerHistogram.lean},
\path{DegreeHistogram.lean} and \path{DegreeMesh.lean} prove
shell bounds, query coverage, clipping, clamping and the $270h$ constant.
\path{DegreeHistogramProbability.lean},
\path{DegreeRateNumbers.lean} and \path{DegreeRelabel.lean} discharge
mesh probabilities, flat branches and relabeling.
\item[Actual-law corollary.]
\path{InDensityClass.fourth_root_actual_law_inverse} in
\path{DegreeLaw.lean} proves Corollary~\ref{cor:law}.
\item[Lower theorem.]
\path{fourth_root_randomized_lower_bound},
\path{fourth_root_lower_exists} and
\path{fourth_root_minimax_obstruction} in \path{LowerBound.lean}
give the finite two-alternative inequality, the existential full-class
form and the $n\ge16$ consequence. Dependencies are
\path{OddGaussian.lean}, \path{LowerAlternatives.lean},
\path{LowerSeparation.lean}, \path{LowerMoments.lean},
\path{LowerTesting.lean}, \path{LowerScale.lean},
\path{LowerProduct.lean} and \path{LowerDataProcessing.lean}.
The success-event hypotheses and independent seed are explicit.
\item[Logarithmic lower theorem and minimax radius.]
\path{logarithmic_randomized_minimax_obstruction} and
\path{logarithmic_lower_bound_at_radius} in \path{LogLowerBound.lean}
prove Theorem~\ref{thm:loglower} and its smaller-radius form.
\path{logarithmic_minimax_radius_bounds} proves
Corollary~\ref{cor:minimax} for each probability seed space, using the
formal infimum \path{orderConfidenceRadius95}.
\path{LogLowerProfile.lean}, \path{LogLowerMoments.lean},
\path{LogLowerAlternatives.lean} and \path{LogLowerPacking.lean}
discharge centered profiles, class membership and separation.
\path{LogLowerTesting.lean}, \path{LogLowerProduct.lean} and
\path{LogLowerScale.lean} discharge testing, iid and seed moments,
and the explicit integer scale. No packing, moment, scale or
good-event premise remains in the final endpoints. The coordinate-data
and all-seed infimum consequences are written consequences of these
arguments; they are not separately named Lean exports.
\end{description}

The proof/evidence revision is
\href{https://github.com/Quantyra/quantyra-null-cone-lean/tree/a0b1f6846d6f0b89a94704ed2954645bcac756fd}{\path{a0b1f6846d6f0b89a94704ed2954645bcac756fd}}.
Lean~4.30.0 and mathlib revision
\path{c5ea00351c28e24afc9f0f84379aa41082b1188f} are pinned.
Authoritative GCP run
\path{space-log-lower-final-acceptance-20261009T171053Z-0c7e20}
completed 3035 root-build jobs and 508 selected exact-type/axiom
audits with zero warnings. These counts cover the library, not
508 results asserted by this paper. All 143 captured proof/check/config
identities and pinned dependencies were verified before and after
execution. The audited axioms are only \path{propext},
\path{Classical.choice} and \path{Quot.sound}. Immutable inputs,
failed attempts and raw diagnostics remain in the repository's
\path{evidence/gcp} tree; its \path{s036-logarithmic-final.json}
indexes final acceptance. Every Lean/Lake invocation ran on GCP.

Manuscript preparation preserves those accepted bytes and invokes no
new Lean build. The companion \path{verify_sources.py} checks the
current and cited-commit hashes against the retained capture and audits;
this byte check does not replace the manual statement-to-prose review.
\path{build.py} compiles this source with Tectonic~0.17.0 and renders
every PDF page with PyMuPDF for inspection. The package records the
PDF hash, page review and compiler diagnostics.
Formal checking establishes the encoded statements, not mathematical
priority, model adequacy or practical efficiency.

\section{Interpretation, practical limits and provenance}
\label{sec:limits}
The result separates existence of a uniformly valid procedure from
useful finite-sample performance. The chosen selector is active only
when $270/\sqrt m<1/2$, requiring $m>291600$ with
$m=\lfloor\sqrt{n/(8\log n)}\rfloor$. The displayed constant is
therefore deliberately conservative. It offers no practical density
resolution at the tested sample sizes.

The separately frozen feasibility study \cite{Study} examined ten
orders across five models at $n=512$ and $3072$, producing thirty
validated report envelopes over three reporting methods. Every density
band had full width one. Twenty exact LP witness pairs documented
the corner limitation. Forty preparation/method runs stayed within
the frozen 600-second and 4-GiB caps; the eighty conditional confirmation
samples were not generated. The degree-filter estimator of this paper
was not implemented or piloted as an additional method. Those results
support neither practical success of this construction nor a
contradiction of its asymptotic guarantee.

At $n=3072$, the lower forbidden radius is approximately $0.0224$,
far below the study's practical target radius $0.45$.
It cannot explain full-width pilot bands or show that the practical
target is impossible. Better constants, efficient observable
reconstruction and improved boundary calibration remain separate
research tasks. The new multi-alternative theorem closes the rate gap,
but its explicit threshold $2^{64}$ and small lower constant were chosen
for a transparent proof, not practical sharpness. The upper constant
650 and the frozen pilot results are unchanged.
No additional order-information penalty or higher-dimensional
reconstruction theorem is proved.

This is a third, separate manuscript family in the shared proof
repository. The reconstruction paper \cite{Reconstruction} addresses
inverse laws and proper time; the gauge paper \cite{Gauge} addresses
an explicit higher-dimensional coordinate obstruction. Their published
artifacts and the frozen study are preserved. This is an unpublished
0.2.0 working revision of the finite-data manuscript. Its published
0.1.0 baseline remains at
\url{https://doi.org/10.5281/zenodo.23265067}; that DOI does not identify
this revised PDF. No publication of the revision is asserted here.

The work was developed with OpenAI Codex assistance under Daniel Eric
Fredriksen's research direction. AI-assisted derivation, formalization
and writing are disclosed. Review follows an open-source,
decentralized and informal workflow: inspectable proofs, reproducible
evidence, downstream use, tests and versioned corrections. Specialist
review and adoption are not prerequisites. Actual feedback must be
distinguished from anticipated uptake. This manuscript and its
companion artifacts are CC-BY-4.0; the proof library is Apache-2.0.

\begin{thebibliography}{99}
\bibitem{Reconstruction} D. E. Fredriksen,
\emph{Quantitative reconstruction from finite causal orders in a conformal
diamond}, version 0.4.0 (2026).
\url{https://doi.org/10.5281/zenodo.23247720}.
\bibitem{Winkler1985} P. Winkler, \emph{Random orders},
Order \textbf{1}, 317--331 (1985).
\url{https://doi.org/10.1007/BF00582738}.
\bibitem{Winkler} P. Winkler, \emph{Random orders of dimension 2},
Order \textbf{7}, 329--339 (PDF dated 1991; publisher metadata 1990).
\url{https://doi.org/10.1007/BF00383197}.
\bibitem{Genest} C. Genest, E. Masiello and K. Tribouley,
\emph{Estimating copula densities through wavelets},
Insurance: Mathematics and Economics \textbf{44}(2), 170--181 (2009).
\url{https://doi.org/10.1016/j.insmatheco.2008.07.006}.
\bibitem{Autin} F. Autin, E. Le Pennec and K. Tribouley,
\emph{Thresholding methods to estimate copula density},
Journal of Multivariate Analysis \textbf{101}, 200--222 (2010).
\url{https://lepennec.perso.math.cnrs.fr/Reprint/Copula/2010-JMVA-ALPT.pdf}.
\bibitem{SeckMamane} C. T. Seck and S. Mamane,
\emph{Strong uniform convergence rates of the linear wavelet estimator
of a multivariate copula density}, South African Statistical Journal
\textbf{58}(1), 35--56 (2024).
\url{https://doi.org/10.37920/sasj.2024.58.1.3}.
\bibitem{Swanepoel} J. Swanepoel, C. T. Seck and S. Mamane,
\emph{A new wavelet estimator of multivariate copula densities based on
Sklar's theorem, with optimal strong uniform convergence rate},
South African Statistical Journal \textbf{58}(2), 99--108 (2024).
\url{https://doi.org/10.37920/sasj.2024.58.2.2}.
\bibitem{Braun} M. Braun, \emph{Spacetime reconstruction by order and
number}, arXiv:2507.01907v1 (2025).
\url{https://arxiv.org/abs/2507.01907v1}.
\bibitem{Gauge} D. E. Fredriksen,
\emph{A certified coordinate-gauge obstruction for causal-order
reconstruction in a $2+1$ diamond}, version 0.1.0 (2026).
\url{https://doi.org/10.5281/zenodo.23249816}.
\bibitem{Hoeffding} W. Hoeffding,
\emph{Probability inequalities for sums of bounded random variables},
Journal of the American Statistical Association \textbf{58}(301),
13--30 (1963).
\url{https://doi.org/10.1080/01621459.1963.10500830}.
\bibitem{Guntuboyina} A. Guntuboyina,
\emph{Lower bounds for the minimax risk using $f$-divergences, and
applications}, arXiv:1002.0042v2 (2011).
\url{https://arxiv.org/abs/1002.0042v2}.
\bibitem{Study} D. E. Fredriksen, \emph{Finite-data feasibility study:
frozen results and decision} (2026), repository revision
\path{0cba197cbef0397209d40ebced29c9dc720573c7}.
\url{https://github.com/Quantyra/quantyra-null-cone-lean/blob/0cba197cbef0397209d40ebced29c9dc720573c7/notes/finite-data-feasibility-results.md}.
\end{thebibliography}
\end{document}
